% GENERATED FILE. Edit canonical lessons, then run npm run generate:books.
% canonical-source-hash: fnv1a64:04f9eec25377419a
% canonical-lessons: MW-C250-manavno, MW-C250-pakvan, MW-C250-uchhav, MW-C250-arti, MW-C250-rolo

\chapter{Lamps and Festive Dishes}
\label{ch:mw-lamps-and-festive-dishes}
\hlchaptermodality{250}

\begin{chapteropening}
\textbf{By the end of this chapter:} \emph{I can say how I celebrated a festival.}

Complete the last lesson of chapter 250: I can say how I celebrated a festival.
\end{chapteropening}

\section[\texorpdfstring{\mw{मनावणो} (manāvṇo)}{मनावणो (manāvṇo)}]{\mw{मनावणो} (manāvṇo) — to celebrate (a festival); to coax}

\label{lesson:MW-C250-manavno}



\begin{quote}
Before the new one: say the Marwadi for Diwali, then the Marwadi for worship.
\end{quote}

\subsection*{You'll want to know: \mw{मनावणो}}
\textbf{\mw{मनावणो}} — \emph{manāvṇo} — \textquotedblleft{}to celebrate (a festival); to coax\textquotedblright{}.

\begin{quote}
\mw{म्हैं} \mw{दिवाळी} \mw{घरै} \mw{मनाई।}
\end{quote}

\begin{quote}
\emph{mhaĩ divāḷī gharai manāī.}
\end{quote}

\begin{quote}
I celebrated Diwali at home.
\end{quote}

Past (I …): \mw{म्हैं} \mw{होळी} \mw{मनाई}, \emph{mhaĩ hoḷī manāī}.

\begin{cousinweb}
\textbf{\mw{मनावणो}} (\emph{manāvṇo}) is inherited from Sanskrit \emph{manyate}, \textquotedblleft{}thinks, honours\textquotedblright{}; CDIAL lists Old Marwari \emph{manāvaï}, \textquotedblleft{}persuades\textquotedblright{}, under it. \textquotedblleft{}To celebrate\textquotedblright{} a feast is probably \textquotedblleft{}to honour\textquotedblright{} it (CDIAL 9857).
\end{cousinweb}

\begin{grammarlens}[title={What you've built}]
One more word for saying what happened.
\end{grammarlens}

\subsection*{Your turn}
\begin{itemize}
\raggedright
  \item \emph{Say it:} \emph{manāvṇo}
  \item \emph{Say it:} \emph{manāvṇo}, once more
  \item \emph{Say it:} \emph{mhaĩ divāḷī gharai manāī.} — I celebrated Diwali at home.
  \item \emph{Recall:} say the Marwadi for Diwali, then the Marwadi for worship, then say \emph{manāvṇo} again
\end{itemize}

\begin{tcolorbox}[breakable,colback=teal!4,colframe=teal!35!black,title={Before you move on}]
What does \mw{मनावणो} mean? (\textquotedblleft{}To celebrate (a festival); to coax\textquotedblright{}.) Say it once more.
\end{tcolorbox}
\section[\texorpdfstring{\mw{पकवान} (pakvān)}{पकवान (pakvān)}]{\mw{पकवान} (pakvān) — festive dishes; fried delicacies}

\label{lesson:MW-C250-pakvan}



\begin{quote}
Before the new one: say the Marwadi for worship, then the Marwadi for to celebrate.
\end{quote}

\subsection*{You'll want to know: \mw{पकवान}}
\textbf{\mw{पकवान}} — \emph{pakvān} — \textquotedblleft{}festive dishes; fried delicacies\textquotedblright{}.

\begin{quote}
\mw{आज} \mw{मां} \mw{पकवान} \mw{बणावै} \mw{है।}
\end{quote}

\begin{quote}
\emph{āj mā̃ pakvān baṇāvai hai.}
\end{quote}

\begin{quote}
Today my mother is making festive dishes.
\end{quote}

A masculine noun.

\begin{cousinweb}
\textbf{\mw{पकवान}} (\emph{pakvān}) goes back to Sanskrit \emph{pakvānna}, \textquotedblleft{}cooked or dressed food\textquotedblright{} (Monier-Williams), built from \emph{pakva}, \textquotedblleft{}cooked\textquotedblright{}, and \emph{anna}, \textquotedblleft{}food\textquotedblright{}.
\end{cousinweb}

\begin{grammarlens}[title={What you've built}]
One more word for saying what happened.
\end{grammarlens}

\subsection*{Your turn}
\begin{itemize}
\raggedright
  \item \emph{Say it:} \emph{pakvān}
  \item \emph{Say it:} \emph{pakvān}, once more
  \item \emph{Say it:} \emph{āj mā̃ pakvān baṇāvai hai.} — Today my mother is making festive dishes.
  \item \emph{Recall:} say the Marwadi for worship, then the Marwadi for to celebrate, then say \emph{pakvān} again
\end{itemize}

\begin{tcolorbox}[breakable,colback=teal!4,colframe=teal!35!black,title={Before you move on}]
What does \mw{पकवान} mean? (\textquotedblleft{}Festive dishes; fried delicacies\textquotedblright{}.) Say it once more.
\end{tcolorbox}
\section[\texorpdfstring{\mw{उछाव} (uchhāv)}{उछाव (uchhāv)}]{\mw{उछाव} (uchhāv) — festive joy; a celebration}

\label{lesson:MW-C250-uchhav}



\begin{quote}
Before the new one: say the Marwadi for to celebrate, then the Marwadi for festive dishes.
\end{quote}

\subsection*{You'll want to know: \mw{उछाव}}
\textbf{\mw{उछाव}} — \emph{uchhāv} — \textquotedblleft{}festive joy; a celebration\textquotedblright{}.

\begin{quote}
\mw{गांव} \mw{मांय} \mw{घणो} \mw{उछाव} \mw{है।}
\end{quote}

\begin{quote}
\emph{gā̃v mā̃y ghaṇo uchhāv hai.}
\end{quote}

\begin{quote}
There is great festive joy in the village.
\end{quote}

A masculine noun.

\begin{cousinweb}
\textbf{\mw{उछाव}} (\emph{uchhāv}) is inherited from Sanskrit \emph{utsāha}, \textquotedblleft{}power, joy\textquotedblright{}; CDIAL lists Old Marwari \emph{uchāha}, \textquotedblleft{}merriment, festival\textquotedblright{}, under it (CDIAL 1882).
\end{cousinweb}

\begin{grammarlens}[title={What you've built}]
One more word for saying what happened.
\end{grammarlens}

\subsection*{Your turn}
\begin{itemize}
\raggedright
  \item \emph{Say it:} \emph{uchhāv}
  \item \emph{Say it:} \emph{uchhāv}, once more
  \item \emph{Say it:} \emph{gā̃v mā̃y ghaṇo uchhāv hai.} — There is great festive joy in the village.
  \item \emph{Recall:} say the Marwadi for to celebrate, then the Marwadi for festive dishes, then say \emph{uchhāv} again
\end{itemize}

\begin{tcolorbox}[breakable,colback=teal!4,colframe=teal!35!black,title={Before you move on}]
What does \mw{उछाव} mean? (\textquotedblleft{}Festive joy; a celebration\textquotedblright{}.) Say it once more.
\end{tcolorbox}
\section[\texorpdfstring{\mw{आरती} (ārtī)}{आरती (ārtī)}]{\mw{आरती} (ārtī) — the waving of a lamp in worship}

\label{lesson:MW-C250-arti}



\begin{quote}
Before the new one: say the Marwadi for festive dishes, then the Marwadi for festive joy.
\end{quote}

\subsection*{You'll want to know: \mw{आरती}}
\textbf{\mw{आरती}} — \emph{ārtī} — \textquotedblleft{}the waving of a lamp in worship\textquotedblright{}.

\begin{quote}
\mw{मां} \mw{आथण} \mw{आरती} \mw{करै} \mw{है।}
\end{quote}

\begin{quote}
\emph{mā̃ āthaṇ ārtī karai hai.}
\end{quote}

\begin{quote}
Mother waves the lamp in worship in the evening.
\end{quote}

A feminine noun.

\begin{cousinweb}
\textbf{\mw{आरती}} (\emph{ārtī}) is inherited from Sanskrit \emph{ārātrika}, the waving of a lamp before an image at night. CDIAL gives no Marwari form in this entry; it lists Hindi and Gujarati \emph{ārtī} (CDIAL 1315).
\end{cousinweb}

\begin{grammarlens}[title={What you've built}]
One more word for saying what happened.
\end{grammarlens}

\subsection*{Your turn}
\begin{itemize}
\raggedright
  \item \emph{Say it:} \emph{ārtī}
  \item \emph{Say it:} \emph{ārtī}, once more
  \item \emph{Say it:} \emph{mā̃ āthaṇ ārtī karai hai.} — Mother waves the lamp in worship in the evening.
  \item \emph{Recall:} say the Marwadi for festive dishes, then the Marwadi for festive joy, then say \emph{ārtī} again
\end{itemize}

\begin{tcolorbox}[breakable,colback=teal!4,colframe=teal!35!black,title={Before you move on}]
What does \mw{आरती} mean? (\textquotedblleft{}The waving of a lamp in worship\textquotedblright{}.) Say it once more.
\end{tcolorbox}
\section[\texorpdfstring{\mw{रोळो} (roḷo)}{रोळो (roḷo)}]{\mw{रोळो} (roḷo) — an uproar; a din}

\label{lesson:MW-C250-rolo}



\begin{quote}
Before the new one: say the Marwadi for festive joy, then the Marwadi for the waving of a lamp in worship.
\end{quote}

\subsection*{You'll want to know: \mw{रोळो}}
\textbf{\mw{रोळो}} — \emph{roḷo} — \textquotedblleft{}an uproar; a din\textquotedblright{}.

\begin{quote}
\mw{गळी} \mw{मांय} \mw{घणो} \mw{रोळो} \mw{है।}
\end{quote}

\begin{quote}
\emph{gaḷī mā̃y ghaṇo roḷo hai.}
\end{quote}

\begin{quote}
There is a great din in the lane.
\end{quote}

A masculine noun.

\begin{cousinweb}
\textbf{\mw{रोळो}} (\emph{roḷo}) is inherited from Sanskrit \emph{rava}, \textquotedblleft{}a roar, a cry\textquotedblright{}; CDIAL lists Marwari \emph{roḷo}, \textquotedblleft{}tumult\textquotedblright{}, under it (CDIAL 10641).
\end{cousinweb}

\begin{grammarlens}[title={What you've built}]
One more word for saying what happened.
\end{grammarlens}

\subsection*{Your turn}
\begin{itemize}
\raggedright
  \item \emph{Say it:} \emph{roḷo}
  \item \emph{Say it:} \emph{roḷo}, once more
  \item \emph{Say it:} \emph{gaḷī mā̃y ghaṇo roḷo hai.} — There is a great din in the lane.
  \item \emph{Recall:} say the Marwadi for festive joy, then the Marwadi for the waving of a lamp in worship, then say \emph{roḷo} again
\end{itemize}

\begin{tcolorbox}[breakable,colback=teal!4,colframe=teal!35!black,title={Before you move on}]
What does \mw{रोळो} mean? (\textquotedblleft{}An uproar; a din\textquotedblright{}.) Say it once more.
\end{tcolorbox}
